%======================================================================
\chapter{Thin-walled members and real materials}\label{ch:tw}
%======================================================================
% Outline drafted 23 Nov 2026. Labels of this chapter: *:tw-*  (bibitem keys tw-*).
% Status and decisions: ch05_thinwalled/README.md. Proposed title (author to confirm):
% "Rods II: thin-walled members and real materials".

\framepart{Outline}
Chapter~\ref{ch:rd} ended with two promises. Euler's load assumes that the cross-section
does not twist and that the material stays elastic. Real steel members break both
assumptions. An open section made of thin walls twists far more easily than it bends, and
it may buckle by twisting at a load well below Euler's. A stocky column yields before it
buckles, and the yield starts early, because rolling and welding leave stresses locked in
the section. Each effect lowers the load; Euler's hyperbola becomes an upper bound.

The thread of the chapter is one structure: the roof of the Hartford Civic Center, which fell
under snow in January 1978. Its top chords were four steel angles bolted into a cross, a
\emph{cruciform}. A cruciform has a remarkable property: its torsional buckling load does not
depend on its length. Bracing it at shorter intervals, the usual cure for a weak strut, does
not help. Its torsional buckling stress also lies close to the yield stress, so that elastic
theory alone cannot tell how much it carries. The chapter follows this column from the
elastic twist to the plastic range, and ends with the design curves that engineers use for
every steel column.

\paragraph{By the end of this chapter you should be able to:}
\begin{itemize}
  \item explain why a thin-walled open section twists before it bends, from the scaling of
        the torsional and bending stiffnesses, and say which model the asymptotic analysis
        of the thin body justifies;
  \item compute the torsional and flexural--torsional buckling loads of a column of open
        section, and the lateral--torsional buckling moment of a beam;
  \item show that the cruciform buckles at $\sigma = G(t/b)^2$ whatever its length, and relate
        this to the buckling of a long plate with a free edge;
  \item compare the tangent-modulus and reduced-modulus loads, and resolve Shanley's
        paradox with a two-degree-of-freedom model;
  \item describe the residual stresses of rolled and welded sections and their effect on the
        stiffness of a column;
  \item derive a design column curve from an imperfect, residually stressed, inelastic
        column, and use the curves of the codes;
  \item compute a column curve with finite elements, residual stresses included, and check it
        against the closed forms.
\end{itemize}

\paragraph{Plan of the chapter.}
Section~\ref{sec:tw-3d} goes from the three-dimensional body to the thin-walled bar: warping,
the torsion constants and the energy of Vlasov. Section~\ref{sec:tw-ft} computes the torsional and
flexural--torsional loads of columns and the lateral--torsional load of beams.
Section~\ref{sec:tw-cruciform} studies the cruciform and the Hartford collapse.
Section~\ref{sec:tw-shanley} leaves elasticity: tangent and reduced modulus, Shanley's model, and
the plastic buckling paradox of the cruciform. Section~\ref{sec:tw-curves} adds residual stresses
and imperfections and arrives at the column curves of the codes. A summary, exercises with
solutions (page~\pageref{sec:tw-exo}) and the references close the chapter.
\draftnote{Sources: Trahair and Ba\v{z}ant--Cedolin (thin-walled theory), Shanley and Hutchinson
(inelastic columns), the SSRC guide and Brush--Almroth \S2.6 (residual stresses, column curves),
Hibbeler \S13.5--13.7 (tangent modulus, design formulas; numbers and examples, file to upload
again), Loomis et al.\ and Martin--Delatte (Hartford). Detailed plan: README of the chapter.}

%======================================================================
\section{From the three-dimensional body to the thin-walled bar}\label{sec:tw-3d}
\draftnote{About 15\,\% of the chapter. (1)~Picture: a strip of width $b$ and thickness $t$ bends
with $EI \propto b^3t$ but twists with $GJ \propto bt^3$; for an open section $GJ/EI \sim (t/b)^2$, so it
twists first. A closed section (tube) has $J \propto$ (area enclosed)$^2t$/perimeter (Bredt) and does
not. (2)~Vlasov kinematics: twist $\varphi$, sectorial coordinate $\omega$, warping displacement
$-\omega\varphi'$, warping constant $\Gamma$, shear centre $(y_0,z_0)$. (3)~Energy: the fibre strain of
Chapter~\ref{ch:rd} for a fibre at $(y,z)$ of a section that also rotates by $\varphi$ gives the
Wagner term $\tfrac12 P r_0^2\varphi'^2$, with $r_0^2 = (I_y + I_z)/A + y_0^2 + z_0^2$:
$\mathcal W = \tfrac12\int(EI_z v''^2 + EI_y w''^2 + GJ\varphi'^2 + E\Gamma\varphi''^2) -
\tfrac12 P\int(v'^2 + w'^2 + r_0^2\varphi'^2 + 2z_0 v'\varphi' - 2y_0 w'\varphi')$ (signs to check against
Trahair). gbox: table of stiffness scalings. (4)~Why it works: Vlasov's theory as a $\Gamma$-limit
of 3D linear elasticity for cross-sections made of rectangles of sides $\varepsilon$ and
$\varepsilon^2$ (Freddi, Morassi and Paroni 2007), closed sections (Davini, Freddi and Paroni 2014),
asymptotic expansions (Trabucho and Vi\~ano 1996). Where it fails: walls so thin that they buckle
locally (plates, Chapter~\ref{ch:pl}), distortion of the section, short members.}

%======================================================================
\section{Torsional and flexural--torsional buckling}\label{sec:tw-ft}
\draftnote{Engineering core. (1)~Adjacent equilibrium: three coupled equations in $v$, $w$,
$\varphi$; pinned ends with fork supports, sine modes, a $3\times3$ determinant, a cubic in $P$.
(2)~Doubly symmetric section (I, H): three separate loads $P_y$, $P_z$ and
$P_\varphi = (GJ + \pi^2E\Gamma/L^2)/r_0^2$. (3)~Monosymmetric sections (angles, tees, channels):
bending about the axis of symmetry stays alone, the other two couple and give a load below
both. Example in numbers: an equal-leg angle strut and a channel, computed by
\texttt{code/python/ft\_section.py}. (4)~Lateral--torsional buckling of a beam:
$M_{cr} = (\pi/L)\sqrt{EI_zGJ}\,\sqrt{1 + \pi^2E\Gamma/(GJL^2)}$. Failures: the Dee Bridge (1847,
Petroski: lateral instability of the girders); the Marcy pedestrian bridge (2002, a girder
system that buckled torsionally during the pour of the deck). (5)~Post-buckling of the
torsional mode: stable, like a plate (forward to Chapter~\ref{ch:pl}); interaction of local and
global modes (Bushnell \S7, the Augusti model of Chapter~\ref{ch:ld}).
Physical box (table top): a ruler held on its edge and loaded at the tip twists sideways; the
critical load scales as $\sqrt{EI_zGJ}/L^2$.}

%======================================================================
\section{The cruciform column}\label{sec:tw-cruciform}
\draftnote{(1)~Four legs of width $b$ and thickness $t$ meeting at one point: the sectorial
coordinate vanishes, $\Gamma$ is only the tiny warping of the legs themselves, $J = 4bt^3/3$,
$I_y + I_z = 4tb^3/3$. gbox: \emph{a cruciform twists at the same stress whatever its length},
$\sigma_T = G(t/b)^2$. Why it works: the twist stiffness $GJ\varphi'^2$ and the Wagner term
$Pr_0^2\varphi'^2$ both carry $\varphi'^2$, so $L$ cancels. (2)~The same number from a plate: a long
plate pinned on one edge and free on the other buckles at $k = 6(1-\nu)/\pi^2 = 0.426$, the
textbook $0.425$; with the legs' own warping one gets $k = 0.425 + (b/a)^2$ (bridge to
Chapter~\ref{ch:pl}). (3)~Torsion or flexure? The crossover length $L = \pi(b^2/t)\sqrt{E/6G}
\approx 2.07\,b^2/t$ for steel: $2$~m for $b = 100$~mm, $t = 10$~mm; shorter cruciforms twist.
(4)~Box: the Hartford Civic Center (1978), a $91\times110$~m space truss whose top chords were
four-angle cruciforms; the snow load was below the design load; the torsional analysis of
Loomis et al.\ (1980) points to torsional buckling of the chords; the bracing assumed in the
computer model was not provided by the offset joints (Martin and Delatte 2001). Lesson: the
model did not contain the mode (Chapter~\ref{ch:nc}, Section~\ref{sec:nc-trust}).
(5)~For usual proportions, $G(t/b)^2$ is close to the yield stress: the cruciform is the bridge
to Section~\ref{sec:tw-shanley}. Physical box (table top): a cruciform of card, cut shorter and
shorter, buckles at the same load; a tube of the same card does not twist.}

%======================================================================
\section{Inelastic buckling: tangent modulus and Shanley's model}\label{sec:tw-shanley}
\draftnote{(1)~Where Euler's hyperbola meets yield (Section~\ref{sec:rd-slender}). Engesser's
tangent-modulus load $P_t = \pi^2E_tI/L^2$ and the reduced-modulus load of Considère, Engesser
and von Kármán ($E_r = 4EE_t/(\sqrt E + \sqrt{E_t})^2$ for a rectangle); tests fall near $P_t$.
(2)~Shanley's paradox (1946--47): the column can bend at $P_t$ under increasing load, without
unloading. Shanley's model: two rigid bars joined by a cell of two flanges, bilinear material,
two degrees of freedom: the method of Chapter~\ref{ch:ld}. Closed form of the path from $P_t$,
rising towards $P_r$; the imperfect model has a maximum load below $P_t$. gbox: \emph{in the
plastic range bifurcation starts at the tangent-modulus load, and the maximum load lies
between $P_t$ and $P_r$.} Uniqueness (Hill) is cited and left to Chapter~\ref{ch:so}.
(3)~Hutchinson's asymptotics of plastic bifurcation (1974) for the initial slope of the path
(check the exponent on the paper). (4)~The plastic buckling paradox of the cruciform: flow
theory keeps the shear modulus $G$ and predicts $G(t/b)^2$ far above the tests; deformation
theory uses the secant shear modulus and agrees; small imperfections reconcile the two
(Onat and Drucker 1953; Hutchinson and Budiansky 1976; reviews: Becque 2016, Shamass 2020;
all to check). (5)~Hibbeler \S13.5: tangent modulus read on a measured stress--strain curve
(worked example to take from the book).}

%======================================================================
\section{Residual stresses and column curves}\label{sec:tw-curves}
\draftnote{(1)~Origin: cooling after rolling (flange tips cool first and end in compression,
about $0.3\sigma_Y$; tension near the web), welding (tension near $\sigma_Y$ at the welds), cold
forming. Self-equilibrated fields $\sigma_0$: $\int\sigma_0\dd A = 0$, $\int\sigma_0 y\dd A = 0$.
(2)~Effect: the section yields gradually; the bending stiffness is that of the elastic core.
Flange tips yield first, so the weak axis loses $(b_e/b)^3$ and the strong axis $b_e/b$: weak-axis
curves are lower. The CRC parabola (Brush and Almroth \S2.6). (3)~Imperfection, residual stress
and yield together: a maximum load, not a bifurcation. Ayrton--Perry: first yield of a bowed
column (the secant formula of Chapter~\ref{ch:rd} with a bow), with the equivalent imperfection
$\alpha(\bar\lambda - 0.2)$, gives the European curves: $\alpha = 0.13$, $0.21$, $0.34$, $0.49$, $0.76$ for
curves $a_0$, $a$, $b$, $c$, $d$, calibrated on more than 1000 ECCS tests; $\chi = 0.60$ on curve $b$
and $0.54$ on curve $c$ at $\bar\lambda = 1$. SSRC curves, AISC. Hibbeler \S13.6--13.7: design
formulas for steel, aluminium and timber on one plot; eccentric loads and interaction.
gbox: \emph{a column curve is the maximum load of an imperfect, residually stressed,
inelastic column; the loss is largest near $\bar\lambda = 1$, where yield and buckling meet}
(erosion of an optimum, Chapter~\ref{ch:ld}). (4)~Box: the Quebec Bridge (1907): a built-up lower
chord with weak latticing (a shear-soft column, Chapter~\ref{ch:nc}); its bow grew from 19 to
57~mm in the two weeks before the collapse, a Southwell signal (Chapter~\ref{ch:rd}) that was
read too late.}

%======================================================================
\framepart{Summary}
\begin{gbox*}
\begin{itemize}
  \item \textbf{To be written.} \draftnote{take-home messages: twist is cheap for open sections;
        three loads and their coupling; the cruciform's length-free load; tangent modulus and
        Shanley; residual stresses and column curves.}
\end{itemize}
\end{gbox*}

%======================================================================
\framepart{Exercises}\label{sec:tw-exo}
\framesub{Short questions}
\draftnote{Planned: why tubes do not twist; the crossover length of a cruciform; which mode for
an angle strut; tangent or reduced modulus from a given curve; Shanley in one sentence; $\chi$
from curve $b$; why the weak axis has a lower curve; what Hartford teaches about a computer
model.}

\framesub{Problems}
\draftnote{Planned: (1)~the three loads of an H-section and the cubic of a channel; (2)~the
cruciform: $G(t/b)^2$, the plate check $k = 0.425 + (b/a)^2$, the crossover length; (3)~Shanley's
model, perfect and imperfect; (4)~the Lehigh residual stress pattern: elastic core, both axes,
the CRC parabola; (5)~Ayrton--Perry and the European curves, against the secant formula;
(6)~design of a steel, an aluminium and a timber column with the code formulas (Hibbeler);
(7)~lateral--torsional buckling of a ruler and of a girder during a deck pour.}

\framesub{Computer exercises}
\draftnote{Planned: (1)~Python: section properties (shear centre, $J$, $\Gamma$) and the three loads
of any thin-walled open section; (2)~Python: Shanley's model by continuation with the
Chapter~\ref{ch:st} code; (3)~Python: a fibre column with residual stresses and a bow, its
maximum load against the slenderness, our own column curve; (4)~Cast3M: multifibre beams with
residual stresses given as initial fibre stresses, elastic--plastic, bow $L/1000$,
\texttt{pasapas} with large displacements, the column curve checked against (3) and the
European curves; a shell model of the cruciform at three lengths (\texttt{flambage}) and in
the plastic range, where Cast3M's flow theory shows the paradox.}

\begin{chapterbib}{99}
\bibitem{tw-trahair1993} N.~S.~Trahair, \emph{Flexural-Torsional Buckling of Structures}, E \& FN Spon, London, 1993.
\bibitem{tw-shanley1947} F.~R.~Shanley, Inelastic column theory, \emph{Journal of the Aeronautical Sciences} 14 (1947) 261--268.
\bibitem{tw-hutchinson1974} J.~W.~Hutchinson, Plastic buckling, \emph{Advances in Applied Mechanics} 14 (1974) 67--144.
\bibitem{tw-ziemian2010} R.~D.~Ziemian (ed.), \emph{Guide to Stability Design Criteria for Metal Structures}, 6th ed., Wiley, Hoboken, 2010.
\bibitem{tw-bazant1991} Z.~P.~Ba\v{z}ant and L.~Cedolin, \emph{Stability of Structures}, Oxford University Press, 1991.
\bibitem{tw-bushnell1981} D.~Bushnell, Buckling of shells---pitfall for designers, \emph{AIAA Journal} 19 (1981) 1183--1226.
\bibitem{tw-freddi2007} L.~Freddi, A.~Morassi and R.~Paroni, Thin-walled beams: a derivation of Vlassov theory via $\Gamma$-convergence, \emph{Journal of Elasticity} 86 (2007) 263--296.
\bibitem{tw-loomis1980} R.~S.~Loomis, R.~H.~Loomis, R.~W.~Loomis and R.~W.~Loomis, Torsional buckling study of Hartford Coliseum, \emph{Journal of the Structural Division (ASCE)} 106 (1980) 211--231.
\bibitem{tw-martin2001} R.~Martin and N.~Delatte, Another look at Hartford Civic Center Coliseum collapse, \emph{Journal of Performance of Constructed Facilities} 15 (2001) 31--36.
\bibitem{tw-pearson2006} C.~Pearson and N.~Delatte, Collapse of the Quebec Bridge, 1907, \emph{Journal of Performance of Constructed Facilities} 20 (2006) 84--91.
\end{chapterbib}
